\section{Realizing the Factors on \tcompile}
\label{sec:analyzers}

Each factor family of Section~\ref{sec:factors} is realized as an \emph{analyzer}: a sweep that takes seeds
from one source, changes one factor, and reuses the oracle, confirmation, and reporting of
Section~\ref{sec:phase3}. Table~\ref{tab:analyzers} lists them with the scale of one sweep on PyTorch~2.14.
The path-side analyzers are new; the program-side ones extend the metadata-driven sweeps of earlier
differential testers with boundary values and unsupported dtypes.

\begin{table}[t]
\caption{Analyzers: the factor that is varied, the reference the result is compared with, and the scale of one
sweep on PyTorch~2.14 (CPU unless noted). Path-side analyzers vary how the program is compiled; program-side
analyzers vary the input or the operator call.}
\label{tab:analyzers}
\scriptsize
\setlength{\tabcolsep}{4pt}
\begin{tabular}{@{}lp{3.05cm}p{4.9cm}p{4.3cm}@{}}
\toprule
Side & Analyzer & Varied factor & Scale of one sweep \\
\midrule
path & graph-break position & a break after every statement (top level or nested), by \code{graph\_break}, \code{disable}, or \code{print}; three backends; \code{dynamic} & 274 programs, 660--816 insertion points, 6~min per backend \\
path & backend layer, context & \code{eager} / \code{aot\_eager} / \code{inductor}; forwarding dispatch mode (forward and gradients); autocast; 17 compile compositions & 955 differentiable op--dtype pairs; 530 ops $\times$ 2 autocast dtypes; 14 programs $\times$ 25 compositions \\
path & artifact route & \code{export} + serialize + load; AOTInductor package + load; runner destroyed after a run; CUDA-graph output reuse (19 patterns) & 520 ops (CPU), 19 patterns (T4) \\
path & context sequences & one factor changed between two calls of one artifact; return trip; cross-process disk cache & 26 switches, 29 cache cases \\
\midrule
program & boundary values & inf, NaN, extremes, near-singular and near-overflow values in vector-length tensors; \code{dynamic}; multi-precision reference for 44 unary and 23 binary functions & 1{,}203--1{,}345 op--dtype pairs per platform (CPU, T4) \\
program & unsupported dtypes & every dtype the eager kernel rejects & 1{,}345 pairs (CPU), 1{,}303 (T4) \\
program & dtype pairs, scalars & $8{\times}8$ dtype pairs in 4 argument forms; 11 binding forms of a scalar & 17{,}664 cases; 386 ops \\
program & source factors (\scs) & one semantic factor of a compilation site at a time; boundary values from predicates & 29 corpus programs, 114 obligations \\
program & Python corpus & 274 programs of Dynamo-modeled Python features, two calls each & 3 backends $\times$ 2 shape modes \\
program & decomposition, differentiation & Inductor decomposition executed eagerly; backward, complex, forward-mode, \code{vmap}; 11 optimizers & 576 ops, 11{,}911 variants; 327--678 ops per mode \\
\bottomrule
\end{tabular}
\end{table}


\paragraph{Graph-break position.} The corpus programs are rewritten at the syntax-tree level so that
\code{torch.\_dynamo.graph\_break()} follows every statement, either at the top level of the function or also
inside loop, conditional, \code{with}, and \code{try} bodies; nested function definitions are left intact and
closures are re-bound to the original cells so that the rewritten function sees the same state. The rewritten
program is compared with eager execution and with the un-rewritten compiled program in the same process, so
that the report lists the divergences the break \emph{added}. Three variants replace the break by a statement
with the same effect through a different mechanism: a call to a function wrapped in Dynamo's \code{disable}
decorator, a \code{print}, and a Python side effect on external state; the first two force a break with and without an
observable effect, the third produces no break and tests side-effect replay alone.

\paragraph{Backend layer and execution context.} The same seed is executed with \code{backend} set to each layer,
with \code{dynamic=True}, under a \code{TorchDispatchMode} whose \code{\_\_torch\_dispatch\_\_} only forwards
the call (the mode is entered by the compile runtime itself, so a difference here means that the compiled
program computes under a mode the user never wrote), under autocast, and under compositions of
\tcompile that occur in application code (a compiled function calling another, \code{compile(compile(f))},
\code{fullgraph}, a warm-up in \code{no\_grad} before a call with gradients, two backends on one function).
For differentiable operators the sweep also compares input gradients. Every difference is confirmed without
\tcompile where possible, which is how three findings were attributed to the dispatch machinery rather than to
the compiler.

\paragraph{Artifact route.} Each OpInfo sample is exported, serialized and deserialized, and executed; packaged
with AOTInductor, loaded, and executed; and the loaded runner is destroyed while its outputs are alive and
before the OpenMP worker threads of its last kernel have parked. Outputs are compared with eager execution and
with \tcompile so that the report states whether a difference belongs to the route or to the shared lowering.

\paragraph{Program-side sweeps.} The boundary-value analyzer replaces the data of every OpInfo sample by the
edge values of Section~\ref{sec:factors} in tensors of at least the vector width, on CPU and CUDA; a variant
compares elementary functions against a 30-digit reference near singularities and near overflow. The
unsupported-dtype analyzer casts each sample to every dtype the eager kernel rejects and records whether the
compiled program raises the same error, another, or none. The corpus analyzer runs the 274 Python programs twice
each and compares return values, state, and exceptions with CPython; the differentiation analyzers compare
backward, complex, forward-mode, and \code{vmap} traversals; the decomposition analyzer executes Inductor's
decompositions eagerly under a dispatch mode, without a C++ compiler in the loop.

\paragraph{Cost.} Path-side analyzers reuse seeds, so their cost is low: a graph-break sweep over the corpus takes
six minutes per backend, the dispatch-mode sweep covers 955 operator--dtype pairs in under an hour, and the artifact
route covers 520 operators in 45 minutes. Resumable per-operator records make a native crash cost one operator,
not the sweep.
